\chapter{Inference on Biological Systems}
\label{ch:inference}

\begin{gomdraftnotice}
This chapter is under active development and contains only a structural outline
with preliminary material. Full exposition, worked examples, proofs, and exercises
will be added in future editions.
\end{gomdraftnotice}

\epigraph{%
  We cannot open the body and read the parameters.
  We must infer them from the traces left by motion.
}{}

\section{Parameter Identification from Motion Data}

Biological system parameters (muscle strengths, tendon slack lengths, segment inertias)
cannot typically be measured directly in vivo. They must be inferred from observable
motion data using optimization and estimation techniques.

\begin{center}
\begin{tikzpicture}[scale=0.9, >=Stealth, node distance=2cm]
  \node[draw, thick, fill=blue!10, minimum width=2cm] (model) {Model $f(\state, \bm{\theta})$};
  \node[draw, thick, fill=blue!10, minimum width=2cm, right of=model] (predict) {Prediction $\hat{\bm{y}}$};
  \node[draw, thick, fill=green!10, minimum width=2cm, align=center, right of=predict] (compare) {Compare\\to data};
  \node[draw, thick, fill=red!10, minimum width=2cm, below of=compare] (optimize) {Optimize $\bm{\theta}$};

  \node[left, font=\tiny] at (model.west) {Parameters};
  \node[below, font=\tiny] at (predict.south) {Observations};

  \draw[->, thick] (model) -- (predict);
  \draw[->, thick] (predict) -- (compare);
  \draw[->, thick] (compare) -- (optimize);
  \draw[<->, thick] (optimize) -- (model);

  \node[draw, thick, fill=orange!10, minimum width=2cm, below of=model] (prior) {Prior $p(\bm{\theta})$};
  \draw[->, dashed] (prior) -- (optimize) node[midway, left, font=\tiny] {Bayes};
\end{tikzpicture}
\end{center}

\subsection{Maximum Likelihood Estimation}

Given a parameterized model $\dot{\state} = f(\state, \bm{\theta})$ with parameters
$\bm{\theta}$ and noisy observations $\bm{y}_k = h(\state_k) + \bm{v}_k$ at times
$t_k$, the maximum likelihood estimate is:
\begin{equation}
  \hat{\bm{\theta}} = \argmin_{\bm{\theta}} \sum_{k=1}^{N}
  \norm{\bm{y}_k - h(\state_k(\bm{\theta}))}_{R^{-1}}^2,
  \label{eq:ch8:mle}
\end{equation}
where $R$ is the measurement noise covariance.

\subsection{Bayesian Estimation with Physical Constraints}

The Bayesian framework (introduced in Chapter~6) is particularly valuable for
biomechanical parameter estimation because it naturally incorporates prior knowledge
from the literature:

\begin{lstlisting}[language=Python, caption={Bayesian parameter estimation for musculoskeletal model.}]
import numpy as np
from scipy.optimize import minimize

def log_posterior(
    theta: np.ndarray,
    observations: np.ndarray,
    forward_model,
    prior_mean: np.ndarray,
    prior_cov: np.ndarray,
    noise_cov: np.ndarray
) -> float:
    """Compute the log-posterior for parameter estimation.

    Parameters
    ----------
    theta : np.ndarray
        Model parameters.
    observations : np.ndarray, shape (n_obs, n_dims)
        Observed data.
    forward_model : callable
        theta -> predicted_observations.
    prior_mean : np.ndarray
        Prior parameter mean.
    prior_cov : np.ndarray
        Prior parameter covariance.
    noise_cov : np.ndarray
        Observation noise covariance.

    Returns
    -------
    float
        Negative log-posterior (for minimization).
    """
    # Log-likelihood
    predicted = forward_model(theta)
    residuals = observations - predicted
    R_inv = np.linalg.inv(noise_cov)
    log_likelihood = -0.5 * np.sum(residuals @ R_inv * residuals)

    # Log-prior
    prior_residuals = theta - prior_mean
    P_inv = np.linalg.inv(prior_cov)
    log_prior = -0.5 * prior_residuals @ P_inv @ prior_residuals

    return -(log_likelihood + log_prior)


def map_estimate(
    observations: np.ndarray,
    forward_model,
    prior_mean: np.ndarray,
    prior_cov: np.ndarray,
    noise_cov: np.ndarray,
    bounds: list[tuple[float, float]] | None = None
) -> np.ndarray:
    """Compute Maximum A Posteriori (MAP) parameter estimate.

    Parameters
    ----------
    observations : np.ndarray
        Observed data.
    forward_model : callable
        Parameter-to-observation mapping.
    prior_mean : np.ndarray
        Prior mean from literature.
    prior_cov : np.ndarray
        Prior uncertainty.
    noise_cov : np.ndarray
        Measurement noise covariance.
    bounds : list of tuple, optional
        Physical bounds on parameters.

    Returns
    -------
    np.ndarray
        MAP parameter estimate.
    """
    result = minimize(
        log_posterior,
        prior_mean,
        args=(observations, forward_model, prior_mean,
              prior_cov, noise_cov),
        method='L-BFGS-B',
        bounds=bounds
    )
    return result.x
\end{lstlisting}

\section{System Identification for Biological Systems}

System identification techniques from control theory (Volume~I) apply to biomechanical
systems with modifications:

\begin{enumerate}
  \item \textbf{Joint impedance identification}: Perturb the limb and measure the
        response. The stiffness, damping, and inertia of the joint can be estimated
        from the transfer function.
  \item \textbf{Muscle parameter identification}: Optimize muscle parameters (F0,
        optimal length, tendon slack length) to match experimental force or torque data.
  \item \textbf{Neural control identification}: Estimate the feedback gains used by
        the nervous system from perturbation experiments.
\end{enumerate}

\section{Machine Learning Approaches}

Modern ML approaches to biomechanical inference include:
\begin{itemize}
  \item Physics-informed neural networks (PINNs) that enforce equations of motion
  \item Gaussian process regression for joint angle estimation from sparse data
  \item Deep learning for markerless motion capture pose estimation
\end{itemize}

\section*{Exercises}

\begin{enumerate}
  \item \textbf{MAP Estimation.}
        Using the provided code, estimate the optimal fiber length of a biceps muscle
        from simulated isometric force data at 5 joint angles.

  \item \textbf{Identifiability Analysis.}
        For a Hill-type muscle model, analyze which parameters can be uniquely estimated
        from isometric force-angle data alone. Which require dynamic data?

  \item \textbf{(Programming.)} Implement a joint impedance identification experiment:
        simulate a perturbed pendulum and estimate stiffness and damping from the response.
\end{enumerate}
